% Options for packages loaded elsewhere
% Options for packages loaded elsewhere
\PassOptionsToPackage{unicode}{hyperref}
\PassOptionsToPackage{hyphens}{url}
\PassOptionsToPackage{dvipsnames,svgnames,x11names}{xcolor}
%
\documentclass[
  spanish,
  11pt,
  a4paper,
]{article}
\usepackage{xcolor}
\usepackage[margin=2.5cm]{geometry}
\usepackage{amsmath,amssymb}
\setcounter{secnumdepth}{5}
\usepackage{iftex}
\ifPDFTeX
  \usepackage[T1]{fontenc}
  \usepackage[utf8]{inputenc}
  \usepackage{textcomp} % provide euro and other symbols
\else % if luatex or xetex
  \usepackage{unicode-math} % this also loads fontspec
  \defaultfontfeatures{Scale=MatchLowercase}
  \defaultfontfeatures[\rmfamily]{Ligatures=TeX,Scale=1}
\fi
\usepackage{lmodern}
\ifPDFTeX\else
  % xetex/luatex font selection
\fi
% Use upquote if available, for straight quotes in verbatim environments
\IfFileExists{upquote.sty}{\usepackage{upquote}}{}
\IfFileExists{microtype.sty}{% use microtype if available
  \usepackage[]{microtype}
  \UseMicrotypeSet[protrusion]{basicmath} % disable protrusion for tt fonts
}{}
\makeatletter
\@ifundefined{KOMAClassName}{% if non-KOMA class
  \IfFileExists{parskip.sty}{%
    \usepackage{parskip}
  }{% else
    \setlength{\parindent}{0pt}
    \setlength{\parskip}{6pt plus 2pt minus 1pt}}
}{% if KOMA class
  \KOMAoptions{parskip=half}}
\makeatother
% Make \paragraph and \subparagraph free-standing
\makeatletter
\ifx\paragraph\undefined\else
  \let\oldparagraph\paragraph
  \renewcommand{\paragraph}{
    \@ifstar
      \xxxParagraphStar
      \xxxParagraphNoStar
  }
  \newcommand{\xxxParagraphStar}[1]{\oldparagraph*{#1}\mbox{}}
  \newcommand{\xxxParagraphNoStar}[1]{\oldparagraph{#1}\mbox{}}
\fi
\ifx\subparagraph\undefined\else
  \let\oldsubparagraph\subparagraph
  \renewcommand{\subparagraph}{
    \@ifstar
      \xxxSubParagraphStar
      \xxxSubParagraphNoStar
  }
  \newcommand{\xxxSubParagraphStar}[1]{\oldsubparagraph*{#1}\mbox{}}
  \newcommand{\xxxSubParagraphNoStar}[1]{\oldsubparagraph{#1}\mbox{}}
\fi
\makeatother

\usepackage{color}
\usepackage{fancyvrb}
\newcommand{\VerbBar}{|}
\newcommand{\VERB}{\Verb[commandchars=\\\{\}]}
\DefineVerbatimEnvironment{Highlighting}{Verbatim}{commandchars=\\\{\}}
% Add ',fontsize=\small' for more characters per line
\usepackage{framed}
\definecolor{shadecolor}{RGB}{241,243,245}
\newenvironment{Shaded}{\begin{snugshade}}{\end{snugshade}}
\newcommand{\AlertTok}[1]{\textcolor[rgb]{0.68,0.00,0.00}{#1}}
\newcommand{\AnnotationTok}[1]{\textcolor[rgb]{0.37,0.37,0.37}{#1}}
\newcommand{\AttributeTok}[1]{\textcolor[rgb]{0.40,0.45,0.13}{#1}}
\newcommand{\BaseNTok}[1]{\textcolor[rgb]{0.68,0.00,0.00}{#1}}
\newcommand{\BuiltInTok}[1]{\textcolor[rgb]{0.00,0.23,0.31}{#1}}
\newcommand{\CharTok}[1]{\textcolor[rgb]{0.13,0.47,0.30}{#1}}
\newcommand{\CommentTok}[1]{\textcolor[rgb]{0.37,0.37,0.37}{#1}}
\newcommand{\CommentVarTok}[1]{\textcolor[rgb]{0.37,0.37,0.37}{\textit{#1}}}
\newcommand{\ConstantTok}[1]{\textcolor[rgb]{0.56,0.35,0.01}{#1}}
\newcommand{\ControlFlowTok}[1]{\textcolor[rgb]{0.00,0.23,0.31}{\textbf{#1}}}
\newcommand{\DataTypeTok}[1]{\textcolor[rgb]{0.68,0.00,0.00}{#1}}
\newcommand{\DecValTok}[1]{\textcolor[rgb]{0.68,0.00,0.00}{#1}}
\newcommand{\DocumentationTok}[1]{\textcolor[rgb]{0.37,0.37,0.37}{\textit{#1}}}
\newcommand{\ErrorTok}[1]{\textcolor[rgb]{0.68,0.00,0.00}{#1}}
\newcommand{\ExtensionTok}[1]{\textcolor[rgb]{0.00,0.23,0.31}{#1}}
\newcommand{\FloatTok}[1]{\textcolor[rgb]{0.68,0.00,0.00}{#1}}
\newcommand{\FunctionTok}[1]{\textcolor[rgb]{0.28,0.35,0.67}{#1}}
\newcommand{\ImportTok}[1]{\textcolor[rgb]{0.00,0.46,0.62}{#1}}
\newcommand{\InformationTok}[1]{\textcolor[rgb]{0.37,0.37,0.37}{#1}}
\newcommand{\KeywordTok}[1]{\textcolor[rgb]{0.00,0.23,0.31}{\textbf{#1}}}
\newcommand{\NormalTok}[1]{\textcolor[rgb]{0.00,0.23,0.31}{#1}}
\newcommand{\OperatorTok}[1]{\textcolor[rgb]{0.37,0.37,0.37}{#1}}
\newcommand{\OtherTok}[1]{\textcolor[rgb]{0.00,0.23,0.31}{#1}}
\newcommand{\PreprocessorTok}[1]{\textcolor[rgb]{0.68,0.00,0.00}{#1}}
\newcommand{\RegionMarkerTok}[1]{\textcolor[rgb]{0.00,0.23,0.31}{#1}}
\newcommand{\SpecialCharTok}[1]{\textcolor[rgb]{0.37,0.37,0.37}{#1}}
\newcommand{\SpecialStringTok}[1]{\textcolor[rgb]{0.13,0.47,0.30}{#1}}
\newcommand{\StringTok}[1]{\textcolor[rgb]{0.13,0.47,0.30}{#1}}
\newcommand{\VariableTok}[1]{\textcolor[rgb]{0.07,0.07,0.07}{#1}}
\newcommand{\VerbatimStringTok}[1]{\textcolor[rgb]{0.13,0.47,0.30}{#1}}
\newcommand{\WarningTok}[1]{\textcolor[rgb]{0.37,0.37,0.37}{\textit{#1}}}

\usepackage{longtable,booktabs,array}
\newcounter{none} % for unnumbered tables
\usepackage{calc} % for calculating minipage widths
% Correct order of tables after \paragraph or \subparagraph
\usepackage{etoolbox}
\makeatletter
\patchcmd\longtable{\par}{\if@noskipsec\mbox{}\fi\par}{}{}
\makeatother
% Allow footnotes in longtable head/foot
\IfFileExists{footnotehyper.sty}{\usepackage{footnotehyper}}{\usepackage{footnote}}
\makesavenoteenv{longtable}
\usepackage{graphicx}
\makeatletter
\newsavebox\pandoc@box
\newcommand*\pandocbounded[1]{% scales image to fit in text height/width
  \sbox\pandoc@box{#1}%
  \Gscale@div\@tempa{\textheight}{\dimexpr\ht\pandoc@box+\dp\pandoc@box\relax}%
  \Gscale@div\@tempb{\linewidth}{\wd\pandoc@box}%
  \ifdim\@tempb\p@<\@tempa\p@\let\@tempa\@tempb\fi% select the smaller of both
  \ifdim\@tempa\p@<\p@\scalebox{\@tempa}{\usebox\pandoc@box}%
  \else\usebox{\pandoc@box}%
  \fi%
}
% Set default figure placement to htbp
\def\fps@figure{htbp}
\makeatother

\ifLuaTeX
  \usepackage{luacolor}
  \usepackage[soul]{lua-ul}
\else
  \usepackage{soul}
\fi


\ifLuaTeX
\usepackage[bidi=basic,shorthands=off]{babel}
\else
\usepackage[bidi=default,shorthands=off]{babel}
\fi
\ifLuaTeX
  \usepackage{selnolig} % disable illegal ligatures
\fi


\setlength{\emergencystretch}{3em} % prevent overfull lines

\providecommand{\tightlist}{%
  \setlength{\itemsep}{0pt}\setlength{\parskip}{0pt}}



 


\makeatletter
\@ifpackageloaded{caption}{}{\usepackage{caption}}
\AtBeginDocument{%
\ifdefined\contentsname
  \renewcommand*\contentsname{Tabla de contenidos}
\else
  \newcommand\contentsname{Tabla de contenidos}
\fi
\ifdefined\listfigurename
  \renewcommand*\listfigurename{Listado de Figuras}
\else
  \newcommand\listfigurename{Listado de Figuras}
\fi
\ifdefined\listtablename
  \renewcommand*\listtablename{Listado de Tablas}
\else
  \newcommand\listtablename{Listado de Tablas}
\fi
\ifdefined\figurename
  \renewcommand*\figurename{Figura}
\else
  \newcommand\figurename{Figura}
\fi
\ifdefined\tablename
  \renewcommand*\tablename{Tabla}
\else
  \newcommand\tablename{Tabla}
\fi
}
\@ifpackageloaded{float}{}{\usepackage{float}}
\floatstyle{ruled}
\@ifundefined{c@chapter}{\newfloat{codelisting}{h}{lop}}{\newfloat{codelisting}{h}{lop}[chapter]}
\floatname{codelisting}{Listado}
\newcommand*\listoflistings{\listof{codelisting}{Listado de Listados}}
\makeatother
\makeatletter
\makeatother
\makeatletter
\@ifpackageloaded{caption}{}{\usepackage{caption}}
\@ifpackageloaded{subcaption}{}{\usepackage{subcaption}}
\makeatother
\usepackage{bookmark}
\IfFileExists{xurl.sty}{\usepackage{xurl}}{} % add URL line breaks if available
\urlstyle{same}
\hypersetup{
  pdftitle={StudyPilot --- Plataforma Inteligente de Estudio},
  pdfauthor={Marvin Andy Calderón Flores},
  pdflang={es},
  colorlinks=true,
  linkcolor={blue},
  filecolor={Maroon},
  citecolor={Blue},
  urlcolor={Blue},
  pdfcreator={LaTeX via pandoc}}


\title{StudyPilot --- Plataforma Inteligente de Estudio}
\usepackage{etoolbox}
\makeatletter
\providecommand{\subtitle}[1]{% add subtitle to \maketitle
  \apptocmd{\@title}{\par {\large #1 \par}}{}{}
}
\makeatother
\subtitle{Documentación Técnica del Proyecto}
\author{Marvin Andy Calderón Flores}
\date{2026-06-22}
\begin{document}
\maketitle

\renewcommand*\contentsname{Tabla de contenidos}
{
\setcounter{tocdepth}{3}
\tableofcontents
}

\section{Resumen Ejecutivo}\label{resumen-ejecutivo}

\textbf{StudyPilot} es una plataforma web de organización y estudio
inteligente desarrollada con R Shiny, diseñada para estudiantes
universitarios. La aplicación integra inteligencia artificial (Google
Gemini 2.5 Flash) para automatizar la extracción de información de
sílabos y horarios desde PDFs, y ofrece un ecosistema completo:
calendario inteligente con fusión de Google Calendar + horario PDF +
bloques de estudio generados por IA, timer Pomodoro resistente a
background throttling, exámenes de práctica, simulación de promedios, y
un Smart Scheduler que distribuye estudio a lo largo del semestre.

La plataforma opera como una \textbf{Progressive Web App (PWA)
Offline-First}: se instala como app nativa en móviles y escritorio,
mantiene la UI visible sin conexión mediante caché de Service Worker +
localStorage, y persiste la sesión del usuario entre recargas de página
(auto-login). Datos almacenados en MongoDB Atlas con autenticación
propia y contraseñas hasheadas con sodium (scrypt).

\textbf{URL de la aplicación:}
\url{https://studypilot.shinyapps.io/plataforma/}

\textbf{Repositorio:} \url{https://github.com/Mcfcalderon/StudyPilot}

\section{Problemática}\label{problemuxe1tica}

\subsection{Contexto}\label{contexto}

En el entorno académico universitario actual, los estudiantes enfrentan
una carga creciente de responsabilidades: múltiples cursos con
evaluaciones de distintos pesos, trabajos grupales, laboratorios,
prácticas calificadas, y exámenes parciales y finales, todo distribuido
en un ciclo de 16 semanas. La gestión manual de esta información consume
tiempo valioso que podría dedicarse al estudio efectivo.

\subsection{Problemas identificados}\label{problemas-identificados}

\begin{enumerate}
\def\labelenumi{\arabic{enumi}.}
\tightlist
\item
  \textbf{Dispersión de información académica}: Los sílabos, horarios,
  evaluaciones y notas están en documentos separados sin una vista
  unificada.
\item
  \textbf{Cálculo manual de promedios}: Los estudiantes realizan
  cálculos manuales para saber su promedio parcial y cuánto necesitan en
  evaluaciones futuras para aprobar.
\item
  \textbf{Falta de visibilidad temporal}: No existe una herramienta que
  muestre qué evaluaciones se aproximan, cuáles están atrasadas, y cuál
  es la prioridad real.
\item
  \textbf{Ingreso manual tedioso}: Cargar los datos de 6+ cursos al
  inicio de cada ciclo requiere horas de trabajo.
\item
  \textbf{Horarios fragmentados}: El horario del PDF de matrícula,
  Google Calendar y bloques de estudio viven en lugares separados sin
  vista unificada.
\item
  \textbf{Herramientas genéricas}: Las aplicaciones existentes no están
  adaptadas al sistema universitario peruano (escala de 0-20, estructura
  de sílabos UTEC).
\end{enumerate}

\section{Análisis Comparativo}\label{anuxe1lisis-comparativo}

\begin{longtable}[]{@{}
  >{\raggedright\arraybackslash}p{(\linewidth - 8\tabcolsep) * \real{0.2000}}
  >{\raggedright\arraybackslash}p{(\linewidth - 8\tabcolsep) * \real{0.2000}}
  >{\raggedright\arraybackslash}p{(\linewidth - 8\tabcolsep) * \real{0.2000}}
  >{\raggedright\arraybackslash}p{(\linewidth - 8\tabcolsep) * \real{0.2000}}
  >{\raggedright\arraybackslash}p{(\linewidth - 8\tabcolsep) * \real{0.2000}}@{}}
\caption{Comparativa de StudyPilot vs.~competidores
principales}\tabularnewline
\toprule\noalign{}
\begin{minipage}[b]{\linewidth}\raggedright
Característica
\end{minipage} & \begin{minipage}[b]{\linewidth}\raggedright
MyStudyLife
\end{minipage} & \begin{minipage}[b]{\linewidth}\raggedright
Notion
\end{minipage} & \begin{minipage}[b]{\linewidth}\raggedright
PowerPlanner
\end{minipage} & \begin{minipage}[b]{\linewidth}\raggedright
\textbf{StudyPilot}
\end{minipage} \\
\midrule\noalign{}
\endfirsthead
\toprule\noalign{}
\begin{minipage}[b]{\linewidth}\raggedright
Característica
\end{minipage} & \begin{minipage}[b]{\linewidth}\raggedright
MyStudyLife
\end{minipage} & \begin{minipage}[b]{\linewidth}\raggedright
Notion
\end{minipage} & \begin{minipage}[b]{\linewidth}\raggedright
PowerPlanner
\end{minipage} & \begin{minipage}[b]{\linewidth}\raggedright
\textbf{StudyPilot}
\end{minipage} \\
\midrule\noalign{}
\endhead
\bottomrule\noalign{}
\endlastfoot
\textbf{Precio} & Gratis & Freemium & \$1.99 & \textbf{Gratis} \\
\textbf{Extracción IA de sílabos} & No & No & No & \textbf{Sí
(batch)} \\
\textbf{Horario desde PDF} & No & No & No & \textbf{Sí} \\
\textbf{Fusión GCal + PDF + IA} & No & No & No & \textbf{Sí} \\
\textbf{Smart Scheduler semestral} & No & No & No & \textbf{Sí} \\
\textbf{Pomodoro (delta-time)} & Sí & No & No & \textbf{Sí} \\
\textbf{Simulador promedios 0-20} & No & No & Sí (GPA) & \textbf{Sí} \\
\textbf{Exámenes práctica IA} & No & No & No & \textbf{Sí} \\
\textbf{PWA Offline-First} & No & No & No & \textbf{Sí} \\
\textbf{Auto-Login persistente} & Sí & Sí & Sí & \textbf{Sí} \\
\textbf{Drag-to-move calendario} & Sí & Sí & No & \textbf{Sí} \\
\textbf{Adaptado a Perú} & No & No & No & \textbf{Sí} \\
\end{longtable}

\section{Arquitectura del Sistema}\label{arquitectura-del-sistema}

\subsection{Diagrama de arquitectura}\label{diagrama-de-arquitectura}

\begin{verbatim}
file:///C:/Users/Marvin/AppData/Local/Temp/RtmpKOqyjl/file32442d404ea9/widget324452307108.html screenshot completed
\end{verbatim}

\begin{figure}

\centering{

\pandocbounded{\includegraphics[keepaspectratio]{StudyPilot_Documentacion_files/figure-pdf/fig-arquitectura-1.pdf}}

}

\caption{\label{fig-arquitectura}Arquitectura general de StudyPilot (PWA
Offline-First)}

\end{figure}%

\subsection{Stack tecnológico}\label{stack-tecnoluxf3gico}

{\def\LTcaptype{none} % do not increment counter
\begin{longtable}[]{@{}
  >{\raggedright\arraybackslash}p{(\linewidth - 4\tabcolsep) * \real{0.3333}}
  >{\raggedright\arraybackslash}p{(\linewidth - 4\tabcolsep) * \real{0.3333}}
  >{\raggedright\arraybackslash}p{(\linewidth - 4\tabcolsep) * \real{0.3333}}@{}}
\toprule\noalign{}
\begin{minipage}[b]{\linewidth}\raggedright
Componente
\end{minipage} & \begin{minipage}[b]{\linewidth}\raggedright
Tecnología
\end{minipage} & \begin{minipage}[b]{\linewidth}\raggedright
Detalle
\end{minipage} \\
\midrule\noalign{}
\endhead
\bottomrule\noalign{}
\endlastfoot
\textbf{Lenguaje} & R & 4.5+ \\
\textbf{Framework Web} & Shiny + bslib & Bootstrap 5 \\
\textbf{Base de datos} & MongoDB Atlas & Cloud, 12 colecciones \\
\textbf{IA generativa} & Google Gemini 2.5 Flash & via \texttt{ellmer},
6 API keys \\
\textbf{PDF} & \texttt{pdftools} & Extracción de texto \\
\textbf{Seguridad} & \texttt{sodium} + \texttt{digest} & scrypt hashing
+ tokens MD5 \\
\textbf{Interactividad} & \texttt{shinyjs}, JavaScript & Delta-time
timer, drag/drop, CRUD \\
\textbf{PWA} & Service Worker + manifest.json & Offline-First,
instalable \\
\textbf{Offline} & localStorage + sw.js & Cache de calendario +
sesión \\
\textbf{Calendar} & \texttt{lubridate} + custom renderer & Fusión
GCal+PDF, anti-duplicados \\
\textbf{Deployment} & shinyapps.io & studypilot account \\
\textbf{Control versiones} & Git + GitHub & Repositorio privado \\
\end{longtable}
}

\subsection{Estructura de archivos}\label{estructura-de-archivos}

\begin{verbatim}
file:///C:/Users/Marvin/AppData/Local/Temp/RtmpKOqyjl/file324447c54ae/widget324460381ce7.html screenshot completed
\end{verbatim}

\begin{figure}

\centering{

\pandocbounded{\includegraphics[keepaspectratio]{StudyPilot_Documentacion_files/figure-pdf/fig-estructura-1.pdf}}

}

\caption{\label{fig-estructura}Estructura de archivos del proyecto
StudyPilot}

\end{figure}%

\section{Funcionalidades}\label{funcionalidades}

\subsection{1. Dashboard}\label{dashboard}

Vista consolidada del progreso académico del ciclo:

\begin{itemize}
\tightlist
\item
  \textbf{Línea temporal} de 16 semanas con semana actual resaltada
\item
  \textbf{6 KPIs}: Progreso \%, Pendientes, Atrasadas, Alta Prioridad,
  Semanas restantes, Promedio general
\item
  \textbf{Tabla de próximas actividades} con prioridad y días restantes
\item
  Promedio general ponderado por créditos (actualizado en tiempo real)
\end{itemize}

\subsection{2. Pomodoro Timer
(Delta-Time)}\label{pomodoro-timer-delta-time}

Timer de estudio reescrito con \textbf{timestamps absolutos}
(\texttt{Date.now()}) en lugar de \texttt{setInterval}, inmune a
background throttling de navegadores:

\begin{itemize}
\tightlist
\item
  Al iniciar: \texttt{endTimestamp\ =\ Date.now()\ +\ remaining}; en
  cada tick: \texttt{remaining\ =\ endTimestamp\ -\ Date.now()}
\item
  Precisión milimétrica incluso si la app se minimiza o la pestaña va a
  segundo plano
\item
  Duración configurable, pausa/reanudación con \texttt{pausedRemaining}
\item
  Música ambiental (white/brown/pink noise) con Web Audio API
\item
  Notificaciones del sistema operativo (\texttt{Notification\ API})
\item
  Registro de sesiones por curso en MongoDB
\end{itemize}

\subsection{3. Calendario Inteligente (Fusión +
CRUD)}\label{calendario-inteligente-fusiuxf3n-crud}

Motor de fusión que unifica tres fuentes en un solo calendario visual:

\subsubsection{\texorpdfstring{Motor de Fusión
(\texttt{fusionar\_horarios})}{Motor de Fusión (fusionar\_horarios)}}\label{motor-de-fusiuxf3n-fusionar_horarios}

\begin{itemize}
\tightlist
\item
  \textbf{GCal es el Jefe}: Al fusionar PDF + GCal, los eventos de
  Google Calendar tienen prioridad absoluta en hora y color.
\item
  \textbf{Anti-duplicados}: Si un evento de GCal se solapa ≥30 min con
  un bloque del PDF, el bloque PDF se descarta.
\item
  \textbf{Anti-colisión de colores}: Si un color del PDF es idéntico a
  uno de GCal, se reasigna automáticamente desde una paleta alternativa
  (\texttt{alt\_palette}).
\item
  \textbf{Colores dinámicos} (\texttt{aplicar\_colores\_cursos}): Busca
  en la tabla de cursos del usuario → heurística por nombre → paleta
  cíclica automática. Nunca sobreescribe colores de GCal ni de IA.
\end{itemize}

\subsubsection{Smart Scheduler (Visión de
Semestre)}\label{smart-scheduler-visiuxf3n-de-semestre}

\begin{itemize}
\tightlist
\item
  Calcula \texttt{weeks\_remaining} hasta la última evaluación pendiente
\item
  Pasa contexto semestral al LLM para distribución progresiva de estudio
\item
  Usa \texttt{Sys.time()} para recortar white space: no agenda antes de
  la hora actual
\item
  \texttt{obtener\_espacio\_libre()} con \texttt{lubridate::interval()}
  resta TODOS los eventos existentes (GCal + PDF + AI previos + horario
  de sueño)
\item
  Bloques de IA con estilo CSS translúcido/punteado (\texttt{cal-ev-ai})
\end{itemize}

\subsubsection{CRUD Interactivo}\label{crud-interactivo}

\begin{itemize}
\tightlist
\item
  \textbf{Click-to-edit}: Todos los eventos son clickeables (AI, PDF,
  GCal). Modal con título, hora inicio/fin, y selector de color (10
  colores).
\item
  \textbf{Drag-to-move}: Arrastre vertical con snap a 15 minutos. JS
  envía \texttt{cal\_drag\_move} → R actualiza AI blocks o crea
  override.
\item
  \textbf{Eliminar}: Eventos AI se eliminan del reactive; eventos
  GCal/PDF se agregan a \texttt{hidden\_events}.
\item
  \textbf{Persistencia MongoDB}: Overrides
  (\texttt{mg\_cal\_overrides\_set}), hidden events
  (\texttt{mg\_cal\_hidden\_set}), AI blocks
  (\texttt{mg\_ai\_blocks\_set}) --- todo se auto-carga al login.
\end{itemize}

\subsubsection{Rendering Defensivo}\label{rendering-defensivo}

\begin{itemize}
\tightlist
\item
  \texttt{horario\_maestro()} reactive centraliza todo: fusion →
  overrides → hidden → colores → timestamps
\item
  \texttt{tryCatch()} global: si algo falla, muestra alert en vez de
  crash
\item
  \texttt{estandarizar\_timestamps()} limpia \texttt{Z},
  \texttt{+offset}, asegura formato \texttt{\%Y-\%m-\%dT\%H:\%M:\%S}
  puro
\item
  \texttt{HOUR\_H\ =\ 50px} sincronizado con CSS
  \texttt{.cal-hour-label} height
\end{itemize}

\subsection{4. Cursos (Batch Upload)}\label{cursos-batch-upload}

\begin{itemize}
\tightlist
\item
  \textbf{Subida múltiple de sílabos}:
  \texttt{fileInput(...,\ multiple\ =\ TRUE)} con límite de 50 MB
\item
  \textbf{Procesamiento batch}: \texttt{withProgress} +
  \texttt{tryCatch} por archivo --- un PDF roto no detiene el lote
\item
  \texttt{gc()} al final del batch para liberar memoria
\item
  Preview + confirmación para archivo individual; auto-save para batch
\item
  Validación de duplicados: si el curso ya existe, lo reemplaza
\item
  \texttt{guardar\_curso\_extraido()} función compartida (single +
  batch)
\end{itemize}

\subsection{5. Modo Examen}\label{modo-examen}

Preparación para evaluaciones con checklist de temas, barra de progreso
visual persistente en MongoDB, guías de estudio, e integración con
Pomodoro.

\subsection{6. IA / Chat}\label{ia-chat}

Chat flotante con Gemini contextual. Resúmenes automáticos con conceptos
clave, mapa conceptual Mermaid, y preguntas de repaso.

\subsection{7. Examen de Práctica}\label{examen-de-pruxe1ctica}

Hasta 20 preguntas (MC + abiertas) generadas por Gemini basadas en temas
del sílabo, con calificación automática y retroalimentación.

\subsection{8. Notas / Calificaciones}\label{notas-calificaciones}

Input numérico (0-20) con cálculo automático de promedio ponderado:

\[\bar{x} = \frac{\sum_{i=1}^{n} \text{nota}_i \times \text{peso}_i}{\sum_{i=1}^{n} \text{peso}_i}\]

Indicador ``Necesitas X en el Y\% restante para aprobar''. Actividades
formativas (peso 0/NA) se agendan pero no afectan el promedio.

\subsection{9. Actividades}\label{actividades}

Tabla interactiva con filtros, edición via modal, eliminación con
confirmación. Indicador de prioridad basado en proximidad y peso.

\subsection{10. Semanal}\label{semanal}

Vista semanal de actividades con navegación y cards por curso.

\section{PWA y Arquitectura
Offline-First}\label{pwa-y-arquitectura-offline-first}

\subsection{\texorpdfstring{Service Worker (\texttt{sw.js}
v3)}{Service Worker (sw.js v3)}}\label{service-worker-sw.js-v3}

\begin{itemize}
\tightlist
\item
  \textbf{HTML/Navegación}: Network-First --- cachea HTML completo en
  cada carga exitosa; si no hay red, sirve desde caché; si no existe
  caché, muestra página offline estilizada con botón ``Reintentar''.
\item
  \textbf{Assets estáticos} (CSS, JS, fonts): Stale-While-Revalidate ---
  sirve cacheado inmediatamente, actualiza en background.
\item
  \textbf{WebSocket/POST}: Nunca interceptados (Shiny los necesita
  libres).
\end{itemize}

\subsection{Opción Nuclear
Anti-Desconexión}\label{opciuxf3n-nuclear-anti-desconexiuxf3n}

Shiny bloquea la UI al perder conexión con overlays grises. StudyPilot
los elimina completamente:

\begin{itemize}
\tightlist
\item
  \textbf{CSS}: Reglas \texttt{!important} que ocultan
  \texttt{\#shiny-disconnected-overlay}, \texttt{.shiny-notification}, y
  fuerzan \texttt{pointer-events:\ auto} en \texttt{html,\ body}.
\item
  \textbf{JS}: \texttt{nukeShinyDisconnectUI()} elimina elementos del
  DOM. Se ejecuta en \texttt{shiny:disconnected}, \texttt{shiny:error},
  con delays múltiples.
\item
  \textbf{MutationObserver}: Vigila \texttt{document.body}
  permanentemente y elimina cualquier elemento de desconexión que Shiny
  inyecte dinámicamente.
\item
  \textbf{Banner offline}: Barra amarilla fija que deshabilita botones
  de API (IA, Sync GCal, Generar Horario).
\end{itemize}

\subsection{Auto-Login (Persistencia de
Sesión)}\label{auto-login-persistencia-de-sesiuxf3n}

Al recargar la página (F5), el usuario entra directamente sin login:

\begin{enumerate}
\def\labelenumi{\arabic{enumi}.}
\tightlist
\item
  \textbf{Login} → R genera token
  \texttt{digest::digest(user\ +\ salt,\ "md5")} → envía
  \texttt{sendCustomMessage("save\_session")} → JS guarda
  \texttt{sp\_user} + \texttt{sp\_token} en localStorage.
\item
  \textbf{Reload} → \texttt{shiny:connected} se dispara → JS lee
  localStorage → espera 500ms (anti race condition) → envía
  \texttt{Shiny.setInputValue("auto\_login")} → R valida token →
  \texttt{auth\_user()} sin contraseña →
  \texttt{shinyjs::hide("auth\_overlay")}.
\item
  \textbf{Logout} → R envía \texttt{execute\_logout} → JS borra TODO
  localStorage → \texttt{window.location.reload(true)}.
\item
  \textbf{Anti-flicker}: \texttt{login\_form\_container} está hidden por
  defecto. Se muestra solo si no hay sesión guardada.
\item
  \textbf{Tab restore}: El tab activo se guarda en
  \texttt{sp\_active\_tab} y se restaura con
  \texttt{bslib::nav\_select()} al auto-login.
\end{enumerate}

\subsection{localStorage Calendar
Cache}\label{localstorage-calendar-cache}

\begin{itemize}
\tightlist
\item
  R envía \texttt{sendCustomMessage("cache\_calendar",\ events)} cada
  vez que \texttt{horario\_maestro()} se actualiza.
\item
  JS guarda el JSON con timestamp en localStorage.
\item
  En modo offline, \texttt{loadCalendarFromCache()} muestra un aviso con
  la fecha del último cache.
\end{itemize}

\subsection{Botón Instalar PWA}\label{botuxf3n-instalar-pwa}

\begin{itemize}
\tightlist
\item
  \texttt{actionButton("btn\_install\_app")} en el navbar (hidden por
  defecto).
\item
  JS captura \texttt{beforeinstallprompt} → hace visible el botón →
  click ejecuta \texttt{deferredPrompt.prompt()} → diálogo nativo del
  OS.
\item
  Si no hay prompt, muestra instrucciones manuales.
\end{itemize}

\subsection{Responsive CSS}\label{responsive-css}

\begin{itemize}
\tightlist
\item
  \texttt{@media\ (max-width:\ 768px)}: Navbar compacta, calendario con
  labels pequeños, oculta aula, scroll 400px.
\item
  \texttt{@media\ (769-1024px)}: Ajustes tablet.
\item
  \texttt{@media\ (display-mode:\ standalone)}: Padding para safe-area
  del status bar en iOS/Android.
\end{itemize}

\section{Seguridad y Privacidad}\label{seguridad-y-privacidad}

\subsection{Autenticación}\label{autenticaciuxf3n}

\begin{itemize}
\tightlist
\item
  Sistema custom de login/registro con contraseñas hasheadas con
  \textbf{sodium} (scrypt)
\item
  Migración automática de contraseñas legacy al primer login
\item
  Auto-login con tokens \texttt{digest} MD5 almacenados en localStorage
  (no contraseñas)
\item
  Heartbeat cada 15s para mantener sesión activa
\end{itemize}

\subsection{Aislamiento de datos}\label{aislamiento-de-datos}

\begin{itemize}
\tightlist
\item
  \textbf{20+ funciones} \texttt{mg\_*} en \texttt{db\_mongo.R} filtran
  por \texttt{user\_id}
\item
  Cada usuario solo ve sus propios datos (cursos, actividades, notas,
  horarios, overrides, AI blocks)
\end{itemize}

\subsection{Rotación de API Keys}\label{rotaciuxf3n-de-api-keys}

\begin{itemize}
\tightlist
\item
  6 claves de Gemini con failover automático via \texttt{ai\_call()}
\item
  Si una clave recibe HTTP 429, rota a la siguiente
\item
  Capacidad total: \textasciitilde90 RPM, \textasciitilde9,000
  solicitudes/día
\end{itemize}

\subsection{Base de datos}\label{base-de-datos}

MongoDB Atlas con conexión TLS/SSL y connection pooling
(\texttt{.mongo\_cache}).

\section{Limitaciones Actuales}\label{limitaciones-actuales}

\subsection{Técnicas}\label{tuxe9cnicas}

\begin{enumerate}
\def\labelenumi{\arabic{enumi}.}
\tightlist
\item
  \textbf{PDFs de imagen}: \texttt{pdftools} no lee PDFs escaneados ---
  solo texto nativo.
\item
  \textbf{Patrón \texttt{globalenv()}}: Variables compartidas en
  environment global --- problemas en multi-worker.
\item
  \textbf{Timeout en IA}: Llamadas largas pueden exceder el timeout de
  shinyapps.io free tier.
\item
  \textbf{API keys hardcodeadas}: Claves en código fuente (mitigado con
  repo privado).
\end{enumerate}

\subsection{Funcionales}\label{funcionales}

\begin{enumerate}
\def\labelenumi{\arabic{enumi}.}
\setcounter{enumi}{4}
\tightlist
\item
  \textbf{Sin recuperación de contraseña}: No hay flujo ``olvidé mi
  contraseña''.
\item
  \textbf{Extracción no determinística}: La IA puede devolver resultados
  variables para el mismo PDF.
\item
  \textbf{Escala fija}: Calificaciones sobre 20 hardcodeado.
\item
  \textbf{Sin exportar datos}: No hay opción de descargar como
  CSV/Excel.
\item
  \textbf{Sin colaboración}: No se comparten datos entre estudiantes.
\end{enumerate}

\subsection{De UX}\label{de-ux}

\begin{enumerate}
\def\labelenumi{\arabic{enumi}.}
\setcounter{enumi}{9}
\tightlist
\item
  \textbf{Nombres truncados}: Dropdowns muestran nombres cortados.
\item
  \textbf{Sin modo oscuro}: Solo tema claro disponible.
\item
  \textbf{Chat sin historial}: El chat no retiene conversaciones entre
  sesiones.
\end{enumerate}

\section{Trabajo Futuro}\label{trabajo-futuro}

{\def\LTcaptype{none} % do not increment counter
\begin{longtable}[]{@{}
  >{\raggedright\arraybackslash}p{(\linewidth - 4\tabcolsep) * \real{0.3333}}
  >{\raggedright\arraybackslash}p{(\linewidth - 4\tabcolsep) * \real{0.3333}}
  >{\raggedright\arraybackslash}p{(\linewidth - 4\tabcolsep) * \real{0.3333}}@{}}
\toprule\noalign{}
\begin{minipage}[b]{\linewidth}\raggedright
Prioridad
\end{minipage} & \begin{minipage}[b]{\linewidth}\raggedright
Mejora
\end{minipage} & \begin{minipage}[b]{\linewidth}\raggedright
Estado
\end{minipage} \\
\midrule\noalign{}
\endhead
\bottomrule\noalign{}
\endlastfoot
Alta & Variables de entorno para API keys & Pendiente \\
Alta & Recuperación de contraseña via email & Pendiente \\
Alta & Límite de intentos de login & Pendiente \\
Media & Escala de calificaciones configurable (0-20, 0-100, GPA) &
Pendiente \\
Media & Exportar notas como CSV/Excel & Pendiente \\
Media & Dashboard de analytics (gráficos de rendimiento) & Pendiente \\
Media & Modo oscuro & Pendiente \\
Baja & Integración con LMS (Canvas, Moodle) & Pendiente \\
Baja & Spaced repetition integrado & Pendiente \\
Baja & Compartir horario entre estudiantes & Pendiente \\
\st{Alta} & \st{Modo offline con Service Workers} & ✅ Completado \\
\st{Media} & \st{Carga batch de múltiples sílabos} & ✅ Completado \\
\st{Alta} & \st{Timeout de sesión / Auto-login} & ✅ Completado \\
\st{Media} & \st{Timer Pomodoro preciso en background} & ✅
Completado \\
\st{Alta} & \st{Calendario editable (CRUD)} & ✅ Completado \\
\st{Alta} & \st{Smart Scheduler semestral} & ✅ Completado \\
\end{longtable}
}

\section{Historial de Desarrollo}\label{historial-de-desarrollo}

{\def\LTcaptype{none} % do not increment counter
\begin{longtable}[]{@{}
  >{\raggedright\arraybackslash}p{(\linewidth - 4\tabcolsep) * \real{0.3333}}
  >{\raggedright\arraybackslash}p{(\linewidth - 4\tabcolsep) * \real{0.3333}}
  >{\raggedright\arraybackslash}p{(\linewidth - 4\tabcolsep) * \real{0.3333}}@{}}
\toprule\noalign{}
\begin{minipage}[b]{\linewidth}\raggedright
Fase
\end{minipage} & \begin{minipage}[b]{\linewidth}\raggedright
Descripción
\end{minipage} & \begin{minipage}[b]{\linewidth}\raggedright
Deploys
\end{minipage} \\
\midrule\noalign{}
\endhead
\bottomrule\noalign{}
\endlastfoot
Chat 1 & Versión inicial con shinymanager, cursos hardcodeados, MongoDB
básico & \textasciitilde8 \\
Chat 2 & Custom auth, lazy MongoDB init, extracción IA de sílabos, fix
disconnects & \textasciitilde6 \\
Chat 3 & Privacidad por usuario, rotación API keys, horario IA, grades
fix, PWA inicial & \textasciitilde10 \\
\textbf{Chat 4} & \textbf{Calendario fusionado, Smart Scheduler
semestral, PWA Offline-First, Auto-Login, batch upload, Pomodoro
delta-time, CRUD drag/drop, migración a studypilot.shinyapps.io} &
\textbf{\textasciitilde12} \\
\end{longtable}
}

\subsection{Bugs críticos resueltos (Chat
4)}\label{bugs-cruxedticos-resueltos-chat-4}

\begin{itemize}
\tightlist
\item
  \textbf{\texttt{estandarizar\_timestamps} eliminada}: Función borrada
  accidentalmente durante edición, causando crash silencioso en
  \texttt{horario\_maestro()} --- calendario mostraba ``No hay eventos''
\item
  \textbf{\texttt{unused\ argument\ (ai\_blocks)}}: Firma de
  \texttt{obtener\_espacio\_libre} revertida sin el parámetro --- Smart
  Scheduler crasheaba
\item
  \textbf{Color loss}: \texttt{estandarizar\_evento} defaulteaba a
  ``blue'', bloqueando la detección por nombre --- todos los eventos
  eran azules
\item
  \textbf{HOUR\_H mismatch}: 48px en código vs 50px en CSS --- eventos
  se desfasaban progresivamente
\item
  \textbf{Race condition auto-login}: \texttt{setInputValue} llegaba
  antes de que R registrara observers --- resuelto con
  \texttt{setTimeout(500ms)}
\item
  \textbf{SW cache stale}: Service Worker servía \texttt{pomodoro.js}
  viejo --- handlers movidos a HTML inline
\item
  \textbf{Pomodoro selector}:
  \texttt{getElementById(\textquotesingle{}pomo\_duration\textquotesingle{})}
  devolvía wrapper div de Shiny, no el input real
\end{itemize}

\section{Conclusiones}\label{conclusiones}

StudyPilot resuelve la dispersión de información académica y la carga
manual de gestión del ciclo integrando en una sola PWA las
funcionalidades de planificador (Google Calendar), estudio enfocado
(Forest), gestión de notas (PowerPlanner), y material de estudio
(Notion). La v4 añade capacidades de nivel SaaS: calendario inteligente
con fusión de fuentes, persistencia de sesión, y operación
offline-first.

El proyecto demuestra la viabilidad de construir aplicaciones web
modernas usando R (Shiny, bslib, mongolite, ellmer) combinado con
técnicas de frontend avanzadas (Service Workers, localStorage,
MutationObserver, delta-time timers) y servicios cloud (MongoDB Atlas,
Google Gemini, shinyapps.io).

\section{Anexos}\label{anexos}

\subsection{A. Dependencias del
proyecto}\label{a.-dependencias-del-proyecto}

\begin{Shaded}
\begin{Highlighting}[]
\FunctionTok{library}\NormalTok{(shiny)        }\CommentTok{\# Framework web}
\FunctionTok{library}\NormalTok{(bslib)        }\CommentTok{\# Bootstrap 5 UI}
\FunctionTok{library}\NormalTok{(DT)           }\CommentTok{\# Tablas interactivas}
\FunctionTok{library}\NormalTok{(mongolite)    }\CommentTok{\# MongoDB driver}
\FunctionTok{library}\NormalTok{(ellmer)       }\CommentTok{\# Google Gemini API}
\FunctionTok{library}\NormalTok{(pdftools)     }\CommentTok{\# Extracción de texto PDF}
\FunctionTok{library}\NormalTok{(sodium)       }\CommentTok{\# Hashing de contraseñas (scrypt)}
\FunctionTok{library}\NormalTok{(digest)       }\CommentTok{\# Tokens MD5 para auto{-}login}
\FunctionTok{library}\NormalTok{(shinyjs)      }\CommentTok{\# JavaScript desde R}
\FunctionTok{library}\NormalTok{(tidyverse)    }\CommentTok{\# Manipulación de datos}
\FunctionTok{library}\NormalTok{(lubridate)    }\CommentTok{\# Manejo de fechas/intervalos}
\end{Highlighting}
\end{Shaded}

\subsection{B. Esquema MongoDB}\label{b.-esquema-mongodb}

{\def\LTcaptype{none} % do not increment counter
\begin{longtable}[]{@{}
  >{\raggedright\arraybackslash}p{(\linewidth - 4\tabcolsep) * \real{0.3333}}
  >{\raggedright\arraybackslash}p{(\linewidth - 4\tabcolsep) * \real{0.3333}}
  >{\raggedright\arraybackslash}p{(\linewidth - 4\tabcolsep) * \real{0.3333}}@{}}
\toprule\noalign{}
\begin{minipage}[b]{\linewidth}\raggedright
Colección
\end{minipage} & \begin{minipage}[b]{\linewidth}\raggedright
Campos principales
\end{minipage} & \begin{minipage}[b]{\linewidth}\raggedright
Filtro user\_id
\end{minipage} \\
\midrule\noalign{}
\endhead
\bottomrule\noalign{}
\endlastfoot
\texttt{users} & user, password (hash), name, admin & No (global) \\
\texttt{courses} & id, name, short, credits, professor, formula,
eval\_day, color & Sí \\
\texttt{activities} & act\_id, course\_id, type, name, code, date, week,
weight, done, is\_calificada & Sí \\
\texttt{grades} & course\_id, code, grade & Sí \\
\texttt{study\_notes} & note\_id, course\_id, text\_content, source,
created\_at & Sí \\
\texttt{schedules} & dia, hora\_inicio, hora\_fin, curso, codigo, aula &
Sí \\
\texttt{syllabi} & course\_id, filename, content, uploaded\_at & Sí \\
\texttt{exam\_checks} & course\_id, topic\_key, checked & Sí \\
\texttt{pomo\_sessions} & course\_id, duration\_min, completed\_at &
Sí \\
\textbf{\texttt{ai\_blocks}} & summary, start, end, color, is\_ai,
source & Sí \\
\textbf{\texttt{cal\_overrides}} & orig\_title, orig\_start, new\_title,
new\_start\_time, new\_end\_time, new\_color & Sí \\
\textbf{\texttt{cal\_hidden}} & key (summary\textbar start) & Sí \\
\end{longtable}
}

\subsection{\texorpdfstring{C. Funciones de IA
(\texttt{ai\_functions.R})}{C. Funciones de IA (ai\_functions.R)}}\label{c.-funciones-de-ia-ai_functions.r}

{\def\LTcaptype{none} % do not increment counter
\begin{longtable}[]{@{}
  >{\raggedright\arraybackslash}p{(\linewidth - 2\tabcolsep) * \real{0.5000}}
  >{\raggedright\arraybackslash}p{(\linewidth - 2\tabcolsep) * \real{0.5000}}@{}}
\toprule\noalign{}
\begin{minipage}[b]{\linewidth}\raggedright
Función
\end{minipage} & \begin{minipage}[b]{\linewidth}\raggedright
Propósito
\end{minipage} \\
\midrule\noalign{}
\endhead
\bottomrule\noalign{}
\endlastfoot
\texttt{ai\_extract\_syllabus()} & Extrae curso/evaluaciones/temas de
PDF de sílabo \\
\texttt{ai\_extract\_schedule()} & Extrae horario de consolidado de
matrícula \\
\texttt{ai\_generate\_summary()} & Genera resumen + conceptos + mapa
conceptual \\
\texttt{ai\_generate\_questions()} & Genera preguntas MC/abiertas para
examen \\
\texttt{ai\_generate\_study\_guide()} & Genera guía de estudio
completa \\
\texttt{generate\_smart\_schedule\_llm()} & Genera bloques de estudio
con contexto semestral \\
\texttt{calcular\_prioridades\_estudio()} & Calcula prioridades por
proximidad y peso \\
\texttt{ai\_call()} & Wrapper central con rotación de 6 API keys \\
\end{longtable}
}

\subsection{\texorpdfstring{D. Funciones de Calendario
(\texttt{google\_cal.R})}{D. Funciones de Calendario (google\_cal.R)}}\label{d.-funciones-de-calendario-google_cal.r}

{\def\LTcaptype{none} % do not increment counter
\begin{longtable}[]{@{}
  >{\raggedright\arraybackslash}p{(\linewidth - 2\tabcolsep) * \real{0.5000}}
  >{\raggedright\arraybackslash}p{(\linewidth - 2\tabcolsep) * \real{0.5000}}@{}}
\toprule\noalign{}
\begin{minipage}[b]{\linewidth}\raggedright
Función
\end{minipage} & \begin{minipage}[b]{\linewidth}\raggedright
Propósito
\end{minipage} \\
\midrule\noalign{}
\endhead
\bottomrule\noalign{}
\endlastfoot
\texttt{estandarizar\_evento()} & Normaliza cualquier df al schema
estándar (7 columnas) \\
\texttt{estandarizar\_timestamps()} & Limpia start/end: quita Z,
+offset, asegura :SS \\
\texttt{aplicar\_colores\_cursos()} & Asigna colores dinámicos con
anti-colisión GCal \\
\texttt{fusionar\_horarios()} & Cruza PDF + GCal con anti-duplicados
(≥30min overlap) \\
\texttt{pdf\_schedule\_to\_events()} & Convierte horario recurrente a
eventos de una semana \\
\texttt{obtener\_espacio\_libre()} & Calcula huecos libres restando busy
intervals \\
\texttt{gcal\_get\_events()} & Descarga eventos ICS de Google Calendar
público \\
\texttt{parse\_ics\_events()} & Parsea formato ICS a dataframe \\
\end{longtable}
}

\subsection{E. Métricas del proyecto}\label{e.-muxe9tricas-del-proyecto}

{\def\LTcaptype{none} % do not increment counter
\begin{longtable}[]{@{}ll@{}}
\toprule\noalign{}
Métrica & Valor \\
\midrule\noalign{}
\endhead
\bottomrule\noalign{}
\endlastfoot
Líneas de código total & \textasciitilde7,200 \\
Archivos R & 7 \\
Archivos web (CSS/JS/SVG/JSON) & 5 \\
Funciones MongoDB & 20+ \\
Funciones IA & 8 \\
Funciones Calendar & 8 \\
Despliegues a producción & \textasciitilde36 \\
Colecciones MongoDB & 12 \\
API keys Gemini & 6 \\
Pestañas de la app & 10 \\
Sesiones de desarrollo & 4 \\
\end{longtable}
}




\end{document}
